\section{Related Work}
\label{sec:related_work}

Our work sits at the intersection of several active research areas: LLM-based agent systems, multi-agent collaboration, decentralized AI infrastructure, mechanism design for strategic AI systems, and trust/verification in AI-generated content. We survey each area and position our contributions relative to the existing landscape.


\subsection{LLM-Based Agent Systems}

The rapid advancement of LLMs has catalyzed a shift from stateless language models to autonomous agents with perception, memory, tool use, and goal-directed decision-making~\cite{xi2023rise, wang2024survey_agents}.
Foundational frameworks include ReAct~\cite{yao2023react}, which interleaves reasoning traces with environment actions, and Toolformer~\cite{schick2024toolformer}, which enables LLMs to self-supervise tool acquisition.
HuggingGPT~\cite{shen2024hugginggpt} uses an LLM as a controller to orchestrate specialized models across modalities, while ToolLLM~\cite{qin2024toolllm} scales tool use to 16,000+ real-world APIs via depth-first decision tree search.
Embodied agents such as Voyager~\cite{wang2023voyager} demonstrate lifelong learning through code-as-action and automatic curriculum generation.
Recent work on agent infrastructure includes AIOS~\cite{mei2024aios}, which proposes OS-level resource management for concurrent agent execution, and benchmarks such as AgentBench~\cite{liu2024agentbench} that reveal substantial capability gaps between commercial and open-source models on agentic tasks.
Comprehensive surveys~\cite{li2024survey, wang2024survey_agents} map the rapidly expanding landscape.

These systems focus on \textit{individual agent capability}---how a single agent reasons, plans, and uses tools---but do not address the economic coordination, quality verification, or trust mechanisms needed when agents from independent organizations must collaborate.
Our ecosystem provides the missing \textit{coordination layer} that enables trustworthy cross-organizational agent interaction.


\subsection{Multi-Agent LLM Collaboration}

A growing body of work explores multi-agent LLM systems for collaborative task solving.
AutoGen~\cite{wu2023autogen} enables flexible multi-agent conversation with customizable roles and human-in-the-loop capabilities.
MetaGPT~\cite{hong2023metagpt} encodes human-like standard operating procedures into multi-agent coordination, assigning specialized software engineering roles to collaborating agents.
CAMEL~\cite{li2023camel} explores autonomous cooperation through role-playing with inception prompting, while ChatDev~\cite{qian2023communicative} models software development as a multi-agent chat chain.
AgentVerse~\cite{chen2023agentverse} investigates emergent group behaviors in dynamic multi-agent collaboration, and CompeteAI~\cite{zhao2023competeai} studies competitive dynamics among LLM agents in market simulations.

Beyond task-oriented collaboration, recent work has explored multi-agent \textit{deliberation and debate}.
Du et al.~\cite{du2024debate} demonstrate that multi-agent debate significantly improves factual accuracy and reasoning.
GPT-Swarm~\cite{zhuge2024gptswarm} optimizes communication topology among agents, showing that connectivity structure substantially affects collective performance.
ProAgent~\cite{zhang2024proagent} builds agents that proactively anticipate partner behavior through theory-of-mind reasoning.
ChatEval~\cite{chan2024chateval} applies multi-agent debate specifically to evaluation, showing diverse agent personas produce more reliable quality assessments than single-judge approaches.
Zhang et al.~\cite{zhang2024cooperative_embodied} decompose multi-agent cooperation into modular LLM components for embodied settings.

All these frameworks share a common architectural assumption: agents are cooperative and governed by a single orchestrator or operate within a single organization.
None addresses the fundamental challenges of cross-organizational trust, strategic misbehavior, or incentive misalignment that arise when agents are deployed by independent entities with divergent objectives---the core challenges our ecosystem is designed to resolve.
Our work complements these frameworks by providing the economic coordination layer that enables trustworthy collaboration among self-interested agents.


\subsection{Decentralized AI Systems}

Decentralized architectures have redefined digital coordination by enabling trustless value exchange without central authority.
Blockchain platforms like Ethereum~\cite{buterin2014ethereum} introduced programmable smart contracts and token incentives to power autonomous applications, while recent architectures adopt rollups and data availability layers~\cite{celestia2025data} for scalability.

At the intersection of decentralization and AI, several lines of work are relevant.
Bittensor~\cite{bittensor2023} creates a token-incentivized peer-to-peer market for machine intelligence, rewarding models based on their contribution to a shared knowledge pool.
Petals~\cite{borzunov2023petals} enables collaborative decentralized inference by distributing model layers across volunteer nodes.
Yuan et al.~\cite{yuan2024decentralized_inference} provide a comprehensive survey of distributed LLM serving techniques, covering model parallelism, communication optimization, and fault tolerance across decentralized nodes.
The emerging field of Decentralized AI (DeAI)~\cite{ding2024deai} outlines research agendas spanning inference, training, data markets, and governance.

Blockchain-based federated learning~\cite{hu2024blockchain_fl} integrates blockchain with federated learning for verifiable aggregation, tamper-proof model provenance, and incentive mechanisms for heterogeneous participants.

However, these systems focus primarily on model training or inference infrastructure rather than on the higher-level problem of \textit{runtime agent collaboration}---coordinating which agent should serve which task, verifying output quality, and distributing rewards among knowledge producers.
Our work extends decentralized coordination to this knowledge-intensive agent ecosystem layer, where tokens regulate not just access but also validation, service routing, and participation quality.


\subsection{Mechanism Design and Strategic Behavior of LLM Agents}

Economic incentives and mechanism design are central to decentralized system viability.
Web3 protocols embed coordination into infrastructure via programmable tokens---a paradigm known as \textit{tokenomics}---governing token issuance, distribution, and utility~\cite{ito2024cryptoeconomics, abadi2024governance, zuo2023set}.
Information elicitation mechanisms such as prediction markets~\cite{wolfers2004prediction} and peer prediction~\cite{miller2005eliciting} provide theoretical foundations for truthful reporting without ground truth---a problem directly relevant to our PoI mechanism, where validators must be incentivized to report honestly despite the absence of verifiable correct answers.

A rapidly growing literature examines the strategic behavior of LLM agents in economic settings.
Duetting et al.~\cite{duetting2024mechanism} investigate mechanism design for LLM agent interactions, designing truthful mechanisms for settings with private information.
Horton~\cite{horton2023homo} demonstrates that LLMs can simulate economic agents whose behavior aligns with classical game-theoretic predictions, supporting our modeling assumption that agent behavior can be captured by Bayesian game theory.
Brookins and DeBacker~\cite{brookins2023gpt_games} systematically evaluate LLM behavior in canonical strategic-form games, finding near-Nash-equilibrium play in many settings.
Akata et al.~\cite{akata2023repeated_games} evaluate LLMs in iterated games, finding sophisticated cooperation strategies.
Fish et al.~\cite{fish2024collusion} demonstrate that LLM agents can learn to collude in pricing games without explicit coordination, highlighting risks relevant to our non-collusion assumption.
Conitzer et al.~\cite{conitzer2024social_choice} apply social choice theory to AI alignment, proposing preference aggregation mechanisms that resonate with our majority-vote consensus design.

We extend tokenomics and mechanism design beyond financial systems to agent ecosystems, where tokens incentivize knowledge validation and honest service provision, and provide formal game-theoretic guarantees (IC, IR, BNE) tailored to the three-role agent interaction structure.


\subsection{Trust, Verification, and Quality Assessment}

LLM hallucination---generating plausible but incorrect outputs---remains a fundamental challenge~\cite{huang2023survey, ji2023hallucination, zhang2023siren}.
Most mitigation strategies focus on individual model improvements, but a growing line of work tackles reliability through \textit{multi-agent verification}.

LLM-as-judge approaches~\cite{zheng2023judging} use LLMs to evaluate other LLMs' outputs, with MT-Bench showing high agreement with human rankings.
Prometheus~\cite{kim2024prometheus} trains open-source evaluation models, enabling decentralized quality assessment without dependence on proprietary models.
SelfCheckGPT~\cite{manakul2023selfcheckgpt} detects hallucinations by sampling multiple responses and checking consistency, while Cohen et al.~\cite{cohen2023crossexam} use cross-examination between LLMs to detect factual errors.
Self-consistency~\cite{wang2023selfconsistency} and chain-of-verification~\cite{dhuliawala2023chainofverification} leverage multiple reasoning paths to improve reliability.
Self-contradictory hallucination detection~\cite{mundler2024selfcontradictory} shows that contradiction frequency correlates with hallucination rate, motivating multi-perspective verification.
Peer evaluation through discussion~\cite{chan2024chateval} demonstrates that multi-agent deliberation produces more reliable assessments than single-judge approaches.

Our approach tackles reliability at the \textit{ecosystem level}: rather than trusting any single agent or relying on self-verification, we aggregate assessments from multiple architecturally heterogeneous validators under economic incentives that promote truthful evaluation.
This connects multi-agent verification with mechanism design: PoI provides the \textit{economic backbone} that makes cross-verification incentive-compatible, ensuring validators invest genuine effort rather than free-riding on others' assessments.
Kapoor et al.~\cite{kapoor2024agents} argue that rigorous evaluation of AI agent capabilities is essential for building reliable systems---our PoI mechanism provides exactly such rigorous, decentralized, and incentive-compatible evaluation.
